\documentclass[11pt]{article}

\usepackage[margin=1in]{geometry}
\usepackage{booktabs}
\usepackage{amsmath}
\usepackage{amssymb}
\usepackage{hyperref}
\usepackage{cite}
\usepackage{authblk}

\title{Glutamate Neurotransmission from Leptin-Receptor Cells in the Ventral Premammillary Nucleus is Required for Female Sexual Maturation and Fertility}

\author[1]{Anonymous Author}
\affil[1]{Anonymous Institution}

\date{}

\begin{document}

\maketitle

\begin{abstract}
Leptin signalling in the ventral premammillary nucleus (PMv) is a permissive gate for puberty and female fertility, but the neurotransmitter mediating the PMv$\to$reproductive-axis output has not been identified. Here we show that chemogenetic activation of leptin-receptor-expressing (LepRb) PMv neurons triggers luteinizing hormone (LH) secretion, and that ablating \textit{Vglut2} (vesicular glutamate transporter 2) from LepRb cells disrupts sexual maturation and fertility in the female mouse. Unilateral stereotaxic injection of AAV-hM3Dq in the PMv of LepRb-Cre females, followed by intravenous (iv.) injection of clozapine-$N$-oxide (CNO) or clozapine at time $0$, produced an increase in LH within $10$--$20$\,min in about $55\%$ (8/15) of ``PMv-hit'' animals ($0.65$--$3.11$\,ng/mL), with a one-sample $t$-test against zero for the average LH change confirming the increase in the hit-with-LH group ($t_7 = 3.54$, $p = 0.009$), and not in PMv-hit no-LH ($t_6 = 1.37$, $p = 0.22$) or negative controls ($t_6 = 1.00$, $p = 0.36$; $n = 7$). The number of cFos-immunoreactive neurons per section in the PMv differed across the three groups (one-way ANOVA $F_{2,19} = 81.66$, $p < 0.0001$) and correlated with peak LH (Pearson $r = 0.67$, $p = 0.0007$). Clozapine produced LH in $4/4$ hM3Dq-injected animals (one-sample $t$-test $t_3 = 8.79$, $p = 0.003$) and in $2/7$ no-AAV controls. In LepRb$^{\Delta}$Vglut2 females, GnRH immunoreactivity in the Arc was reduced ($t_5 = 8.01$, $p = 0.0005$), the LH response to kisspeptin-10 was attenuated (mixed-effects model, genotype main effect $F_{1,11} = 8.26$, $p = 0.015$; AUC $t_{11} = 3.14$, $p = 0.009$), mediobasal Tac2 expression dropped (about $50\%$; $t_8 = 3.39$, $p = 0.009$), cycle duration was longer ($t_{25} = 3.71$, $p = 0.001$), and $25\%$ of females were sterile over $60$ days paired with proven breeders ($75\%$ produced at least one litter, vs.\ $100\%$ of controls). Rescuing LepR and concurrently deleting Vglut2 in PMv neurons in LepR$^{\mathrm{loxTB}}$;Vglut2$^{\mathrm{flox}}$ females prevented puberty-onset recovery (vaginal opening) and pregnancy (one-way ANOVAs for PMv pSTAT3 $F_{3,13} = 73.53$, $p < 0.0001$; Arc pSTAT3 $F_{3,16} = 171.8$, $p < 0.0001$; corpora lutea $F_{3,17} = 31.16$, $p < 0.0001$). Glutamate from PMv LepRb neurons is thus required for female sexual maturation and fertility.
\end{abstract}

\section{Introduction}

Leptin is a permissive metabolic signal for puberty and adult reproduction \cite{elias2010review,hill2018,manfredilozano2016,cravo2011,kauffman2007kndy}. LepRb neurons in the ventral premammillary nucleus (PMv) are a central hub of this control: deletion of LepRb in PMv neurons delays puberty in mice, and AAV-Cre-mediated re-expression of LepR in PMv of LepR$^{\mathrm{loxTB}}$ mice restores puberty and fertility \cite{donato2011,mahany2018}. About $75\%$ of PMv LepRb neurons are depolarized by leptin \cite{williams2011}, and the proximate gatekeeper of gonadotrophin release, the GnRH neuron, receives strong glutamatergic drive. The identity of the neurotransmitter relaying PMv LepRb drive to the reproductive axis has been unclear, however. Kisspeptin is a direct GnRH stimulator but only about $10\%$ of KNDy (Kiss1) neurons in the arcuate nucleus express LepR \cite{cravo2013,donato2011,louis2011,true2011}, suggesting that leptin acts indirectly through an intermediary. Here we test whether glutamate released from LepRb PMv neurons is the critical relay, combining chemogenetic activation of PMv LepRb neurons with selective Vglut2 deletion from LepRb cells (LepRb$^{\Delta}$Vglut2) and with an endogenous-LepR re-expression model in LepR$^{\mathrm{loxTB}}$;Vglut2$^{\mathrm{flox}}$ females.

\section{Methods}

\paragraph{Chemogenetics.} Adult LepRb-Cre females (8--10 weeks old) received unilateral stereotaxic injections of $\sim 50$\,nL of AAV-hM3Dq (fused to mCherry) in the PMv. Coordinates (from rhinal vein): AP $-5.4$\,mm; ML $\pm 0.5$\,mm (from superior sagittal sinus); DV $-5.4$\,mm (from dura). Mice were connected to a Culex Automated Blood Sampling System $3$ days after surgery. The following day, $3$ baseline blood samples were taken every $10$\,min, followed by an iv.\ injection of CNO ($n=19$) or clozapine ($n=4$) at $0.5$\,mg/kg in $50\,\mu$L at time $0$. Samples were collected every $10$\,min for $1$\,h. Negative controls ($n=7$) did not receive an AAV and were injected with clozapine only. For AUC analysis, the $-10$\,min value served as baseline.

\paragraph{Kisspeptin challenge.} Baseline LH was established before an intraperitoneal kisspeptin-10 injection at $65\,\mu$g/kg \cite{wang2019kiss}, and blood was sampled every $7$\,min for $30$\,min.

\paragraph{LepRb$^{\Delta}$Vglut2 mice.} Conditional deletion of Vglut2 in LepRb cells was verified by chromogenic BaseScope Duplex ISH for Lepr (Green) and Vglut2 (Fast Red) ($n=4$ LepRb-Cre; $n=4$ LepRb$^{\Delta}$Vglut2). Body weight was monitored weekly from $13$ to $26$ weeks. Puberty onset (vaginal opening in females, balanopreputial separation in males), estrous cycles, and fertility (over $60$ days with proven breeders) were tracked.

\paragraph{LepR re-expression with concurrent Vglut2 deletion.} LepR$^{\mathrm{loxTB}}$;Vglut2$^{\mathrm{flox}}$ females ($n=9$) received AAV-Cre in the PMv. Wild-type mice were positive (fertile) controls ($n=13$; $6$ with AAV-GFP, $7$ without). LepR$^{\mathrm{loxTB}}$ (LepR-null) mice were negative (infertile) controls ($n=4$). Six to eight weeks after surgery, animals received an ip.\ leptin injection ($3$\,mg/kg) $1.5$\,h before perfusion, and PMv and Arc pSTAT3 were quantified.

\paragraph{Histology, ISH and qPCR.} Brains were fixed and sectioned at $30\,\mu$m on a freezing microtome; reproductive organs were processed paraffin-embedded with H\&E staining. Single ISH for Vglut2 used a T7-transcribed $^{35}$S-riboprobe hybridized overnight at $57\,^\circ$C. qPCR was performed on $20$\,ng cDNA in triplicate $10\,\mu$L reactions (SYBR Green Taq MasterMix; Bio-Rad CFX-384) for Gnrhr, LH$\beta$, FSH$\beta$, Kiss1, Pdyn, Kiss1r, Tac3r, Tac2, and CGA, normalized to Actb or $\beta$2m. Data are mean $\pm$ SEM. LH functional sensitivity was $0.016$\,ng/mL (range $0.016$--$4.0$\,ng/mL); intra-assay CV $2.2\%$; inter-assay CVs $7.3\%$/$5.0\%$/$6.5\%$ at Low/Medium/High QC.

\section{Results}

\subsection{Chemogenetic activation of PMv LepRb neurons triggers LH release}
Of $15$ mice with unilateral AAV-hM3Dq centred in the PMv (``PMv-hit''), about $55\%$ ($8/15$) showed an increase in LH above $0.5$\,ng/mL within $10$--$20$\,min of CNO ($0.65$--$3.11$\,ng/mL). Two mice had injections outside the PMv (``PMv-miss''; $n=2$); five AAV-mCherry animals served as additional controls. A modest LH increase of $0.57$\,ng/mL was observed in one PMv-miss animal and an increase of $0.93$\,ng/mL at $30$\,min was observed in one AAV-mCherry animal. One-sample $t$-tests on the average LH change at $-10$, $10$ and $20$\,min showed a significant rise in the PMv-hit LH-increase group ($n=8$, $t_7 = 3.54$, $p = 0.009$), but not in PMv-hit no-LH-increase ($n=7$, $t_6 = 1.37$, $p = 0.22$) or negative controls ($n=7$, $t_6 = 1.00$, $p = 0.36$).

\subsection{PMv cFos tracks LH response}
PMv cFos-ir counts per section differed across the three groups (one-way ANOVA $F_{2,19} = 81.66$, $p < 0.0001$), with PMv-hit/LH-increase animals ($n=9$) higher than PMv-hit/no-LH ($n=7$) and negative controls ($n=7$). The correlation between PMv cFos counts and peak LH was $r = 0.67$, $p = 0.0007$ (Pearson).

\subsection{Clozapine reproduces the CNO response}
Clozapine administered iv.\ at time $0$ produced an LH rise of $0.87$--$1.18$\,ng/mL at $10$\,min in all $4$ AAV-hM3Dq animals (one-sample $t$-test $t_3 = 8.79$, $p = 0.003$), versus $2/7$ no-AAV controls ($28.5\%$), which did not reach significance as a group ($t_6 = 1.83$, $p = 0.12$). PMv cFos differed between hM3Dq-clozapine and no-AAV animals ($t_8 = 17.46$, $p < 0.0001$).

\subsection{Posterior Arc contamination is additive, not determinant}
Peak LH correlated with the number of cFos-ir $+$ mCherry-ir neurons in the posterior Arc (Pearson $r = 0.62$, $p = 0.002$; Kruskal-Wallis group test $= 14.63$, $p < 0.0001$). POMC-ir coexpression was observed in $9.3 \pm 2.6\%$ of cFos-ir Arc neurons in PMv-hit animals ($12.3 \pm 5.0$ neurons) and $9.6 \pm 3.1\%$ in hM3Dq-clozapine animals ($2.3 \pm 1.3$ neurons). Of these POMC/cFos-ir neurons, about $50\%$ also coexpressed mCherry in PMv-hit animals ($7.3 \pm 3.2$ neurons) but almost none in hM3Dq-clozapine animals ($0.3 \pm 0.3$ neurons). In LepRb-Cre;Kiss1-hrGFP AAV-hM3Dq-injected PMv-hit animals ($n=4$), only $2.5 \pm 0.95\%$ of cFos-ir Arc neurons were Kiss1-hrGFP$^+$ ($1.8 \pm 0.3$ neurons), suggesting at most $1$--$2$ KNDy neurons were directly activated.

\subsection{LepRb$^{\Delta}$Vglut2 mice develop late-onset obesity}
Adult LepRb$^{\Delta}$Vglut2 mice (13--26 weeks) developed late-onset obesity due to increased fat mass in both sexes (males fat mass $t_9 = 3.77$, $p = 0.004$; females fat mass $t_{11} = 3.7$, $p = 0.0035$). Body-mass genotype main effects: males $F_{1,10} = 13.1$, $p = 0.0047$, with post-hoc week 20 $t_{140} = 3.27$, $p = 0.019$ and week 26 $t_{140} = 5.47$, $p < 0.0001$; females $F_{1,11} = 4.34$, $p = 0.06$, with post-hoc week 23 $t_{154} = 4.05$, $p = 0.0011$, week 25 $t_{154} = 3.1$, $p = 0.03$, and week 26 $t_{154} = 5.59$, $p < 0.0001$.

\subsection{LepRb$^{\Delta}$Vglut2 mice exhibit delayed/subtle puberty}
Vaginal opening age ($t_{18} = 0.84$, $p = 0.41$) and body mass at VO ($t_{18} = 1.47$, $p = 0.16$) did not differ; first estrus was delayed ($t_{14} = 2.45$, $p = 0.028$) at equal body mass ($t_{14} = 0.59$, $p = 0.56$). In males, age at balanopreputial separation (BPS) did not differ ($t_{11} = 1.80$, $p = 0.1$) but body mass at BPS did ($t_{11} = 2.43$, $p = 0.033$). Adult female cycle duration was longer ($t_{25} = 3.71$, $p = 0.001$). Time in estrus trended longer ($t_{26} = 1.97$, $p = 0.059$); proestrus was longer ($t_{26} = 2.81$, $p = 0.009$) and diestrus was longer ($t_{26} = 2.39$, $p = 0.024$). Number of cycles in 30 days was lower (LepRb-Cre $4.47 \pm 0.30$ vs.\ LepRb$^{\Delta}$Vglut2 $3.31 \pm 0.36$; $t_{26} = 2.45$, $p = 0.021$).

\subsection{LepRb$^{\Delta}$Vglut2 females show subfertility}
$75\%$ of LepRb$^{\Delta}$Vglut2 vs.\ $100\%$ of control mice produced at least one litter over $60$ days paired with proven breeders. Days from mating to first litter did not differ (Mann-Whitney, $p = 0.32$). Among fertile females, first-litter size did not differ (LepRb-Cre $7.86 \pm 1.14$ vs.\ LepRb$^{\Delta}$Vglut2 $7.50 \pm 1.41$; $t_{11} = 0.20$, $p = 0.85$).

\subsection{Reduced Arc GnRH, Tac2, and LH response to kisspeptin}
GnRH-ir integrated density in the Arc was reduced in LepRb$^{\Delta}$Vglut2 ($t_5 = 8.01$, $p = 0.0005$; $n = 4$ Vglut2$^{\mathrm{flox}}$ vs.\ $n = 3$ LepRb$^{\Delta}$Vglut2). Basal LH was not different ($t_{10} = 1.91$, $p = 0.08$). After iv.\ kisspeptin-10 ($65\,\mu$g/kg), LH levels were higher in controls than in floxed females (mixed-effects: time $F_{1.14,12.19} = 15.91$, $p = 0.001$; genotype $F_{1,11} = 8.26$, $p = 0.015$; time$\times$genotype $F_{3,32} = 5.00$, $p = 0.006$; Sidak post-hoc at $7$\,min $t_{10.97} = 3.79$, $p = 0.01$; $14$\,min $t_{9.4} = 3.22$, $p = 0.04$). Post-injection AUC differed ($t_{11} = 3.14$, $p = 0.009$). Pituitary Gnrhr ($t_{10} = 0.86$, $p = 0.41$), LH$\beta$ ($t_{10} = 1.47$, $p = 0.17$), FSH$\beta$ ($t_9 = 1.63$, $p = 0.14$) and CGA ($100 \pm 17.47\%$ vs.\ $245.3 \pm 89.91\%$; $t_8 = 1.20$, $p = 0.26$) were similar. In the mediobasal hypothalamus (relative to B2m), Kiss1 ($t_8 = 1.42$, $p = 0.19$), Pdyn ($t_8 = 0.02$, $p = 0.98$), Kiss1r ($t_8 = 1.44$, $p = 0.19$), Tac3r ($t_8 = 0.99$, $p = 0.35$), and Tac2 ($t_8 = 3.39$, $p = 0.009$) were measured; only Tac2 was significantly reduced, by about $50\%$.

\subsection{Endogenous LepR re-expression with Vglut2 deletion blocks puberty}
We targeted LepR$^{\mathrm{loxTB}}$;Vglut2$^{\mathrm{flox}}$ females ($n=9$) with AAV-Cre in the PMv and compared them with wild-type ($n=13$), LepR$^{\mathrm{loxTB}}$ (LepR-null, $n=4$), and AAV-Cre-injected LepR$^{\mathrm{loxTB}}$ ($n=5$) females. Successful PMv-targeted LepR$^{\mathrm{loxTB}}$;Vglut2$^{\mathrm{flox}}$ injections ($n=6$) showed increased leptin-induced pSTAT3-ir in PMv (one-way ANOVA across four groups: $F_{3,13} = 73.53$, $p < 0.0001$), and Vglut2 mRNA was decreased relative to AAV-Cre LepR$^{\mathrm{loxTB}}$ ($t_8 = 4.20$, $p = 0.003$). Arc pSTAT3 also differed across groups ($F_{3,16} = 171.8$, $p < 0.0001$) and the number of corpora lutea across groups ($F_{3,17} = 31.16$, $p < 0.0001$). Correctly PMv-injected LepR$^{\mathrm{loxTB}}$;Vglut2$^{\mathrm{flox}}$ females showed no vaginal opening and no pregnancy to the end of the experiment ($\approx 60$ days).

\begin{table}[h]
\centering
\caption{Key outcomes from the four-group PMv manipulation in the LepR$^{\mathrm{loxTB}}$/Vglut2$^{\mathrm{flox}}$ rescue model.}
\label{tab:rescue}
\begin{tabular}{lc}
\toprule
Comparison & Test statistic \\
\midrule
PMv pSTAT3-ir (4 groups) & $F_{3,13} = 73.53$, $p < 0.0001$ \\
Arc pSTAT3-ir (4 groups) & $F_{3,16} = 171.8$, $p < 0.0001$ \\
Corpora lutea (4 groups) & $F_{3,17} = 31.16$, $p < 0.0001$ \\
PMv Vglut2 mRNA (AAV-Cre groups) & $t_8 = 4.20$, $p = 0.003$ \\
\bottomrule
\end{tabular}
\end{table}

\section{Discussion}

Together, the four lines of evidence presented here argue that glutamate released by LepRb-expressing PMv neurons is required for female sexual maturation and adult reproduction. Chemogenetic activation of PMv LepRb neurons rapidly triggers LH release ($0.65$--$3.11$\,ng/mL within $10$--$20$\,min) in $55\%$ of PMv-hit animals, tightly correlates with PMv cFos ($r = 0.67$, $p = 0.0007$), and is reproduced by clozapine. LepRb$^{\Delta}$Vglut2 females show delayed first estrus ($t_{14} = 2.45$, $p = 0.028$), longer estrous cycles ($t_{25} = 3.71$, $p = 0.001$), reduced number of cycles in 30 days ($t_{26} = 2.45$, $p = 0.021$), and $25\%$ sterility over $60$ days, despite unaltered pituitary LH$\beta$, FSH$\beta$, Gnrhr and CGA and largely unaltered mediobasal hypothalamic Kiss1, Pdyn, Kiss1r, and Tac3r. The specific $\sim 50\%$ reduction of Tac2 ($t_8 = 3.39$, $p = 0.009$), together with reduced Arc GnRH-ir ($t_5 = 8.01$, $p = 0.0005$) and the attenuated LH response to kisspeptin-10 (genotype $F_{1,11} = 8.26$, $p = 0.015$; AUC $t_{11} = 3.14$, $p = 0.009$), pinpoint a deficit in the neurokinin B/kisspeptin arm of the KNDy network upstream or at the level of GnRH, rather than a pituitary defect. Finally, LepR re-expression with concurrent Vglut2 deletion in PMv neurons of LepR$^{\mathrm{loxTB}}$;Vglut2$^{\mathrm{flox}}$ females blocks puberty onset and pregnancy even though leptin-induced PMv pSTAT3 and Arc pSTAT3 are restored, showing that glutamate release (not LepR signalling per se) is the necessary node.

Our interpretation is that the PMv glutamate signal is necessary, not sufficient. Contamination of POMC \cite{manfredilozano2016} and KNDy \cite{cravo2013,donato2011,louis2011,true2011} neurons in the posterior Arc contributed to the quantitative variability of the CNO-triggered LH response --- peak LH correlated with mCherry/cFos double-labelling in Arc ($r = 0.62$, $p = 0.002$) --- but POMC contamination is modest ($\approx 9\%$ of cFos neurons, $\approx 50\%$ of which mCherry$^+$), and only $1$--$2$ KNDy neurons were directly activated. The consistent dissociation between PMv-hit/LH-increase animals and PMv-miss or AAV-mCherry controls, together with the cFos--LH correlation and the clozapine reproducibility, indicate that the LH response is PMv-driven and not a spontaneous pulsatile event.

Limitations include the modest sample size in the PMv-miss and control groups, the unilateral injection strategy (which likely underestimates the maximum achievable LH release), and the subtle phenotype of LepRb$^{\Delta}$Vglut2 in male reproduction. Future work should test whether the glutamatergic PMv output converges on Tac2 neurons via monosynaptic connections, whether stimulation of glutamatergic PMv axons to the arcuate is sufficient to trigger LH pulses, and whether alternative metabolic signals \cite{hill2010,hill2018} can substitute for leptin in the absence of PMv glutamate. Our data place glutamate from LepRb PMv neurons at the neurochemical heart of the leptin-metabolic control of fertility.

\bibliographystyle{unsrt}
\bibliography{references}

\end{document}
